\subsection{Prescribed finite graph gauges}\label{sec:gauges}

The reference-family construction used the radial, conical, and compact
coefficient maps.  The lateral paths additionally change between radius
and height variables and pass through long height charts near a collapsing
circular profile.  The estimates below retain each highest second
derivative in its bulk norm while controlling its coefficient by
first-derivative traces.

All complex geometric expressions below use the complex bilinear extension of
Euclidean products.  The parameters of a surface chart remain real.  A complex
contour is specified in its reference functions; it does not mean that an
unknown smooth function has been continued to a complex spatial argument.
Square roots of area elements are continued from the positive real branch.
The compact observation charts and their cutoffs remain fixed.

For an interval $I$ of length between two fixed positive constants, set
\begin{align}
 \|v\|_{\mathcal X_h(E;I)}={}&
 \|(v,v_t,v_{\theta\theta},E v_{tt},\sqrt E v_{t\theta})
                         \|_{L^2(I;H^h(\mathbb S^1))}
 \notag\\[-2pt]
 &+\|(v,v_\theta,\sqrt E v_t)\|_{L^\infty(I;H^h(\mathbb S^1))},
 \qquad
 \|f\|_{\mathcal Y_h(I)}=\|f\|_{L^2(I;H^h(\mathbb S^1))}.
 \label{gauge:norm}
\end{align}
Here and below $h\ge4$ is fixed.  The trace terms are supplied by
Lemma~\ref{lin:strip}; one can equivalently include them in the norm.
On overlapping intervals the weights $E$ are comparable, with bounded
normalized derivatives.  A uniformly local weighted norm is the supremum
of the strip norms after division by the chosen growth weight.  All estimates
in this subsection hold at any fixed larger angular order as well.

\begin{lemma}[Mixed prescribed gauges]\label{gauge:mixed}
Let $B\ge1$ and, on a bounded interval, prescribe smooth circular reference
functions and transverse coefficients
\begin{equation}
 X(t,\theta)=(B\zeta(t,\theta),R(t,\theta)e_r(\theta)),\quad
 \zeta=\zeta_0+c_1u,\qquad R=R_0+c_2u.
 \label{gauge:immersion}
\end{equation}
Assume that $R_0$, $\zeta_0'$ and
$c_2\zeta_0'-c_1R_0'$ are bounded away from zero, that their required
fixed derivatives are bounded, and that the circular reference has the
strict ellipticity margins of Lemma~\ref{lin:strip} after division by its
drift.  Constants may depend on these fixed margins.
Then the scalar stationary operator divided by the normal component of
$X_u$ is an analytic map from a sufficiently small
$\mathcal X_h(B^{-2})$ ball to $\mathcal Y_h$.
If $\|u\|_{\mathcal X_h(B^{-2})}\le W$ and $BW$ is sufficiently small,
its nonlinear remainder satisfies
\begin{equation}
 \|\mathcal Q(u)-\mathcal Q(0)-D\mathcal Q(0)u\|_{\mathcal Y_h}
       \le C_h(BW^2+B^2W^3).
 \label{gauge:mixed-bound}
\end{equation}
On that ball its Lipschitz constant, after subtraction of the linear part,
is at most $C_h(BW+B^2W^2)$.  Averaging the equation in these prescribed
coordinates gives the same bound for the defect of the circular mean curve.
These statements include all angular commutations through order $h$.
\end{lemma}

\begin{proof}
Write $a=\zeta_t$, $b=R_t$ and
\begin{equation}
 D=c_2a-c_1b
  =c_2\zeta_0'-c_1R_0'+(c_2c_1'-c_1c_2')u.
 \label{gauge:transverse}
\end{equation}
In particular $D$ contains neither $u_t$ nor $u_\theta$.
The tangent vectors give
\begin{align}
 g_{tt}&=B^2a^2+b^2=:A,
 &g_{t\theta}&=(B^2ac_1+bc_2)u_\theta,
 &g_{\theta\theta}&=R^2+(B^2c_1^2+c_2^2)u_\theta^2,\notag\\
 J:=\det g&=R^2A+B^2D^2u_\theta^2,
 &g^{tt}&=\frac{R^2+(B^2c_1^2+c_2^2)u_\theta^2}{J},\notag\\
 g^{t\theta}&=-\frac{(B^2ac_1+bc_2)u_\theta}{J},
 &g^{\theta\theta}&=\frac A J.
 \label{gauge:metric}
\end{align}
The cancellation in the determinant is exact.  In particular, a term
$B^2\zeta_\theta^2$ cannot be estimated independently of the off-diagonal
square in the determinant.  For $E=B^{-2}$, $d=a^2+Eb^2$,
$L=c_1a+Ec_2b$ and $T=R^2d+D^2u_\theta^2$, the same identities become
\begin{equation}
 g^{tt}=\frac E d+\frac{L^2u_\theta^2}{dT},\qquad
 g^{t\theta}=-\frac{Lu_\theta}{T},\qquad
 g^{\theta\theta}=\frac dT,\qquad
 g^{tt}-\frac{(g^{t\theta})^2}{g^{\theta\theta}}=\frac E d.
 \label{gauge:schur}
\end{equation}
This is the exact Schur complement identity.

An unnormalized normal, sufficient for taking normal components, is
\begin{equation}
 N=\left(-b/B,\ a e_r+
          \frac{b\zeta_\theta-aR_\theta}{R}e_\theta\right).
 \label{gauge:normal}
\end{equation}
It satisfies $N\cdot X_t=N\cdot X_\theta=0$ and $N\cdot X_u=D$.
Consequently the stationary equation is precisely
\begin{equation}
 \mathcal Q(u)=g^{ij}\frac{N\cdot X_{ij}}D
                         +\frac{N\cdot X}{2D}=0.
 \label{gauge:operator}
\end{equation}
No connection term remains in $N\cdot X_{ij}$.  Direct differentiation yields
\begin{align}
 \frac{N\cdot X_{tt}}D&=u_{tt}+C_{tt},
 &C_{tt}&=\frac{a(R_0''+c_2''u+2c_2'u_t)
                    -b(\zeta_0''+c_1''u+2c_1'u_t)}D,\notag\\
 \frac{N\cdot X_{t\theta}}D&=u_{t\theta}+C_{t\theta}u_\theta,
 &C_{t\theta}&=\frac{ac_2'-bc_1'}D-\frac bR,\notag\\
 \frac{N\cdot X_{\theta\theta}}D
        &=u_{\theta\theta}-\frac{aR}D-\frac{2c_2u_\theta^2}R,
 &\frac{N\cdot X}{D}&=\frac{Ra-\zeta b}D.
 \label{gauge:numerators}
\end{align}
Thus the complete operator is affine in the second derivatives.
$C_{tt}$ is at most quadratic in $u_t$ and $C_{t\theta}$ at most affine
in $u_t$, with coefficients analytic in $u$.  Moreover
\begin{align}
 Ra-\zeta b={}&R_0\zeta_0'-\zeta_0R_0'
 +(c_2\zeta_0'+R_0c_1'-c_1R_0'-\zeta_0c_2')u\notag\\
 &+(R_0c_1-\zeta_0c_2)u_t+(c_2c_1'-c_1c_2')u^2.
 \label{gauge:support}
\end{align}
The terms $u u_t$ cancel.  The normalized support function is affine in
$u_t$.  At circular data, irrespective of its speed,
\begin{equation}
       g^{tt}=A^{-1}=O(B^{-2}),\qquad
       g^{t\theta}=0,\qquad g^{\theta\theta}=R^{-2}.
 \label{gauge:circular-cancellation}
\end{equation}
In particular the circular angular curvature term is $-a/(DR)$,
which is affine in $u_t$; it has no unsuppressed quadratic speed term.

Here is a full coefficient estimate.  Use the Banach algebra
$\mathcal A=L^\infty(I;H^h(\mathbb S^1))$ and its module
$\mathcal Y_h$.  By \eqref{gauge:norm},
\begin{equation}
 \begin{gathered}
 \|(u,u_\theta)\|_{\mathcal A}\le W,\qquad
 \|u_t\|_{\mathcal A}\le BW,\\
 \|(u_t,u_{\theta\theta})\|_{\mathcal Y_h}\le W,\qquad
 \|u_{t\theta}\|_{\mathcal Y_h}\le BW,\qquad
 \|u_{tt}\|_{\mathcal Y_h}\le B^2W.
 \end{gathered}
 \label{gauge:slots}
\end{equation}
The functions $R,D,d,T$ are analytic units in this algebra when $BW$
is small.  Equations \eqref{gauge:schur}--\eqref{gauge:numerators}
partition the entire operator into the following classes:
\begin{enumerate}
\item The circular radial coefficient is $E$ times an analytic function
of $(u,u_t)$.  Its change multiplies $u_{tt}$ with bound
$C E(W+BW)B^2W\le CBW^2$.
Its Taylor terms on bounded reference derivatives satisfy this bound too.
\item Every extra radial coefficient has two factors $u_\theta$, by
\eqref{gauge:schur}.  Its product with $u_{tt}$ is bounded by $CB^2W^3$.
\item The mixed coefficient has a factor $u_\theta$ and analytic dependence
on $(u,u_t,u_\theta)$.  Its product with $u_{t\theta}$ is $O(BW^2)$;
any additional speed factor costs at most $BW$.
\item The angular coefficient equals $R^{-2}$ plus a term divisible by
$u_\theta^2$.  Its change multiplying $u_{\theta\theta}$ is
$O(W^2)+O(W^3)$, with analytic speed factors uniformly bounded.
\item The remaining terms are exactly the lower terms in
\eqref{gauge:numerators} and \eqref{gauge:support}.
Pure speed squares in $C_{tt}$ retain the factor $E$; all other pure
speed dependence at $u_\theta=0$ is affine.  Terms with angular tilt
retain its one or two explicit factors.  A product containing only
first derivatives is estimated with one factor in $\mathcal Y_h$,
so even $u_t^2$ costs $BW^2$, rather than $B^2W^2$.
\end{enumerate}
This list exhausts the three inverse-metric contractions and the support
term.  Every further Taylor factor is small in $\mathcal A$, being
bounded by $C(W+BW)$.  The convergent analytic series therefore sum to
\eqref{gauge:mixed-bound}.  Differencing a product gives the asserted
Lipschitz estimate.  The multiplication inequality
$\|fg\|_{H^h}\le C_h\|f\|_{H^h}\|g\|_{H^h}$ proves these statements
with all $h$ angular commutations, keeping exactly one second derivative
in $L^2$; none is put into a trace.

For the last assertion let $U=\langle u\rangle_\theta$ and
$\mathcal N(u)=\mathcal Q(u)-\mathcal Q(0)-D\mathcal Q(0)u$.
The linearized coefficients are independent of $\theta$, whence
\begin{equation}
 \mathcal Q(U)=\mathcal N(U)-\langle\mathcal N(u)\rangle_\theta
 \quad\hbox{if }\mathcal Q(u)=0.
 \label{gauge:averaging}
\end{equation}
The circular curve is exactly $(B(\zeta_0+c_1U),R_0+c_2U)$.
If instead one expands about $U$, its second derivatives need not have
uniform traces: split them into reference derivatives and perturbation
derivatives and keep the latter in the same $L^2$ slot.  The coefficient
Taylor remainder then costs at most $B^2W^3$.  This proves the claim also
for a varying circular mean.
\end{proof}

\begin{lemma}[Long polar height charts]\label{gauge:height}
Prescribe the polar height chart
\[
 X=(Zg(r,\theta),r e_r(\theta))
\]
on real parameters along a fixed complex $r$ contour.  Let $|Z|\asymp B$,
$t=-\log r$, and assume a uniformly bounded circular reference $g_0$ satisfies
$|(g_0)_t|\asymp |r|^2$, with its fixed log derivatives of this size.
Assume $|Br|\ge Y_0$ and put $E=|Br|^{-2}$ on each strip.  The reference
conditions include the full real pullback and drift division.  More precisely,
if $t=t(\xi)$ on the real outward parameter and $J=t_\xi$, assume
$J,J^{-1}$ and their required fixed derivatives are bounded.  Let
$\mathcal F_t(g)$ denote the left side of the log-normalized equation
\eqref{gauge:polar-normalized} below, including its graph-dependent
$H_1$ division.  Write its circular linearization in the log coordinate as
\[
 D\mathcal F_t(g_0)v=(1+d_{\mathrm{ref}})v_t-a_{\mathrm{ref}}v_{tt}
       -c_{\mathrm{ref}}v_{\theta\theta}+q_{\mathrm{ref}}v.
\]
Its full pulled drift is
$\Delta_{\mathrm{ref}}=1+d_{\mathrm{ref}}
             +a_{\mathrm{ref}}J_\xi/J^2$.
Assume this is a bounded unit and that, after division by it,
\[
 \operatorname{Re}\frac{a_{\mathrm{ref}}}{J\Delta_{\mathrm{ref}}}\ge cE,
 \quad \left|\frac{a_{\mathrm{ref}}}{J\Delta_{\mathrm{ref}}}\right|\le CE,
 \quad \operatorname{Re}\frac{Jc_{\mathrm{ref}}}{\Delta_{\mathrm{ref}}}\ge c,
 \quad \left|\frac{Jc_{\mathrm{ref}}}{\Delta_{\mathrm{ref}}}\right|\le C.
\]
The normalized coefficient derivatives and lower terms satisfy the bounds
of Lemma~\ref{lin:strip}.  All metric and graph denominators in
$\mathcal F_t$ remain bounded analytic units in their normalized scales
on the graph neighborhood.  These are algebraic hypotheses
on the prescribed reference, checked for the collapsing paths below.
Let $\operatorname{pb}_J$ mean the algebraic pullback obtained by replacing
$\partial_t$ with $J^{-1}\partial_\xi$ and reading all prescribed
coefficients along the real path.  We use the fixed circular row
\begin{equation}\label{gauge:fixed-row}
 \mathcal F_{\mathrm{norm}}(g)
   =\Delta_{\mathrm{ref}}^{-1}J\,\operatorname{pb}_J\mathcal F_t(g).
\end{equation}
In particular $\Delta_{\mathrm{ref}}$ is held fixed when this map is
linearized at the possibly nonexact profile $g_0$.  This circular row
depends only on $\xi$, so it preserves mean and oscillatory projections.
All strip norms use
the real parameter $\xi$ after this pullback.  Take $Y_0$ large.  The map
$\mathcal F_{\mathrm{norm}}$ is analytic whenever
\begin{equation}
 \frac{E^{-1/2}}{|r|^2}(W_m+W_o)\ll1,
 \qquad
 \frac{E^{-1/2}}{|r|^2}W_o\ll1,
 \label{gauge:height-small}
\end{equation}
where $v=g-g_0$ has mean and oscillation of
$\mathcal X_h(E)$ sizes $W_m,W_o$ respectively.
Writing $W=W_m+W_o$, its remainder
$\mathcal N(v)=\mathcal F_{\mathrm{norm}}(g_0+v)
-\mathcal F_{\mathrm{norm}}(g_0)-D\mathcal F_{\mathrm{norm}}(g_0)v$ obeys
\begin{equation}
 \|\mathcal N(v)\|_{\mathcal Y_h}
 \le C_h\left(\frac{E^{-1/2}}{|r|^2}W^2
              +\frac{E^{-1}}{|r|^4}W_o^2W\right).
 \label{gauge:height-bound}
\end{equation}
With $A_r=E^{-1/2}/|r|^2$ and $C_r=E^{-1}/|r|^4$, its mean and
oscillatory outputs satisfy, respectively,
\begin{align}
 \|\langle\mathcal N(v)\rangle\|_{L^2_\xi}
       &\le C_h\{A_r(W_m^2+W_o^2)+C_rW_o^2W\},\notag\\
 \|\mathcal N(v)-\langle\mathcal N(v)\rangle\|_{\mathcal Y_h}
       &\le C_h\{A_r W_oW+C_rW_o^2W\}.
 \label{gauge:height-projections}
\end{align}
The corresponding multilinear difference estimates hold.  These estimates
are uniform at each fixed angular order.
\end{lemma}
\begin{proof}
The physical slopes in the polar orthonormal frame are
$p=-Zg_t/r$, $q=Zg_\theta/r$.  With $H=1+p^2+q^2$, the curvature
matrix has entries
\begin{equation}
 A=\frac{1+q^2}{H},\qquad D_\theta=\frac{1+p^2}{H},\qquad
 C=-\frac{pq}{H}.
 \label{gauge:polar-matrix}
\end{equation}
The stationary equation, multiplied by $r^2/Z$, is exactly
\begin{equation}
 A(g_{tt}+g_t)+D_\theta(-g_t+g_{\theta\theta})
 -2C(g_{t\theta}+g_\theta)
       =-\tfrac12r^2(g_t+g).
 \label{gauge:polar-equation}
\end{equation}
Complex tangent factors on a prescribed contour are included by the chain
rule in \eqref{gauge:fixed-row}.  Its row is the fixed circular unit,
including the $J_\xi$ term; the graph-dependent $H_1$ normalization remains
inside $\mathcal F_t$ and its linearization.  The preceding normalized
inequalities retain the required strict margins.

Set $p_0=-Z(g_0)_t/r$, $e=p_0^{-2}$, $s=(p-p_0)/p_0$ and $q=p_0 v_1$.
Then $|e|\asymp E$, $s= v_t/(g_0)_t$ and
$v_1=-v_\theta/(g_0)_t$, and the exact coefficient formulas read
\begin{equation}
 A=\frac{e+v_1^2}{e+(1+s)^2+v_1^2},\quad
 D_\theta=\frac{e+(1+s)^2}{e+(1+s)^2+v_1^2},\quad
 C=-\frac{(1+s)v_1}{e+(1+s)^2+v_1^2}.
 \label{gauge:relative-coefficients}
\end{equation}
For an explicit check of the drift division, write instead
$P=g_t/r^2$, $Q=g_\theta/r^2$, $e_1=(Zr)^{-2}$ and
\[
 H_1=(1-r^2/2)P^2-(1+r^2/2)Q^2-r^2e_1/2.
\]
Equation~\eqref{gauge:polar-equation} becomes
\begin{equation}
 g_t-a_1g_{tt}-b_1g_{\theta\theta}
       +2c_1^*g_{t\theta}+2c_1^*g_\theta-d_1g=0,
 \label{gauge:polar-normalized}
\end{equation}
where
\[
 a_1=\frac{e_1+Q^2}{H_1},\qquad
 b_1=\frac{e_1+P^2}{H_1},\qquad
 c_1^*=\frac{PQ}{H_1},\qquad
 d_1=\frac{r^2(e_1+P^2+Q^2)}{2H_1}.
\]
Here $P$ is a bounded unit at the reference.  Put
$\mathfrak p=E^{-1/2}W/|r|^2$ and
$\mathfrak q=E^{-1/2}W_o/|r|^2$.  In the angular algebra these exact
normalized coefficients satisfy
\begin{equation}
 \frac{\|a_1-a_{1,0}\|}{E}
 +\frac{\|b_1-b_{1,0}\|}{E}
 +\frac{\|d_1-d_{1,0}\|}{|r|^2E}
 \le C(\mathfrak p+\mathfrak q^2),\qquad
 \frac{\|c_1^*\|}{\sqrt E}\le C\mathfrak q.
 \label{gauge:relative-normalized}
\end{equation}
In particular at zero tilt the speed variation of $b_1$ retains a
factor $E$, and that of $d_1$ a factor $r^2E$; their remaining variations
have two tilt factors.  Drift division introduces no unweighted
radial-speed term.

Consequently $A-A_0=O(E s)+O(v_1^2)$,
$D_\theta-1=-v_1^2/(e+(1+s)^2+v_1^2)=O(v_1^2)$, and
$C=O(v_1)$, as analytic algebra estimates.
Every pure radial-speed variation in $A$ retains the factor $E$;
angular changes of $A$ have two tilt factors.  The condition that the
mixed coefficient be small relative to $\sqrt E$ is precisely the
second smallness condition in \eqref{gauge:height-small}.

For the full map, assign the traces of $v,v_\theta$ size $W$, the trace
of $v_t$ size $E^{-1/2}W$, and the bulk sizes of
$v_{tt},v_{t\theta},v_{\theta\theta}$ respectively
$E^{-1}W,E^{-1/2}W,W$.  A pure speed change in $A$ times $v_{tt}$ is
bounded by
$E|r|^{-2}(E^{-1/2}W)(E^{-1}W)=A_rW^2$.
A mixed tilt times $v_{t\theta}$ costs $A_rW_oW$.
The two-tilt part of $A$ times $v_{tt}$ costs $C_rW_o^2W$.
The angular coefficient costs less.  These are all second derivative
terms in \eqref{gauge:polar-equation}.
For a coefficient Taylor remainder on a background derivative, use its
size $O(|r|^2)$.  The pure speed part has the factor $E$, and use one
$v_t$ in bulk rather than trace.  For instance
$E|r|^{-4}v_t^2 O(|r|^2)$ has bound
$E^{1/2}|r|^{-2}W^2\le A_rW^2$.
Tilt squares have bound $|r|^{-2}W_o^2\le A_rW_o^2$.
For the normalized zeroth-order term $d_1g_0$, use the factor $r^2E$
in \eqref{gauge:relative-normalized} and $g_0=O(1)$, giving these same
bounds; it is not necessary to assign $g_0$ the derivative size $O(r^2)$.
All lower terms are already displayed in \eqref{gauge:polar-equation};
the support term before division is affine in $g,g_t$.
The analytic $H_1$ division already included in $\mathcal F_t$ preserves
these bounds by its unit expansion.  The fixed row in
\eqref{gauge:fixed-row} and its bounded chain-rule coefficients preserve
them as well.  Higher powers are absorbed using
\eqref{gauge:height-small}.

Finally rotation of $\theta$ commutes with the equation.  At a circular
input, its linearization has circular coefficients.  A mean output
therefore has no term linear in the oscillation, and an oscillatory output
has no term containing only mean inputs.  Applying the preceding
multilinear estimates with these absent terms gives
\eqref{gauge:height-projections}.  Analytic composition in $H^h$ and the
single highest-derivative bulk slot prove every angular commutation and
every difference estimate asserted above.
\end{proof}
